\section{Virtual Machine and Container}
To consistently replicate the software environment needed for debugging the TKJHAT platform across diverse computing infrastructures, we implemented two host-side environments: Virtual Machines and Containers. In the context of our overall debugging pipeline (detailed in Figure \ref{fig:general_structure}), these software environments represent the setup of the host computer, which resides at the left-most boundary of the system.

To ensure toolchain parity, both environments are based on the Debian Linux family. The container uses a lightweight Debian 12 (Bookworm) image to minimise memory footprint and boot times \cite{rosen_namespaces_slides}. Conversely, the Virtual Machine utilises Debian 13 (Trixie) to support modern desktop packages and hardware drivers for the graphical interface. Both maintain strict operational consistency despite these version differences.

A functional debugging setup requires a software stack running on top of the operating system. The software stack used can be divided into four categories:

\begin{description}
    \item[Compiler Toolchain:] The suite of softwares needed to compile C/C++ source code into machine-readable binaries (e.g., UF2, ELF). The compile toolchain's components used for this thesis are \texttt{gcc-arm-none-eabi} cross-compilation infrastructure, which includes \texttt{libnewlib-arm-none-eabi} and \texttt{libnewlib-dev}, used as they provide the standard C library implementation, including memory management, mathematical functions, and system call stubs required for embedded I/O operations. To support the compilation infrastructure, \texttt{CMake} serves as the build system, and \texttt{Ninja} as the generator. Together, they provide a versatile and reliable framework for the Computer System course project and for more advanced embedded systems projects.
    
    \item[Middleware and SDKs:] To unlock all the features supported by the Raspberry Pi Pico platform and complex task management features, the environments also came prepared with Pico and FreeRTOS SDK. These provide the essential hardware abstraction and real-time scheduling capabilities required to support the course's embedded systems projects.
    
    \item[Development Interface:] To provide a familiar but versatile interface for project development, VS Code was selected for its ecosystem and flexibility, serving as the standardised IDE for the course. Alongside the main IDE, helper extensions are installed to provide effective interaction between the IDE and the development platform:
    
    \begin{itemize}
        \item \textbf{Cortex-Debug:}\label{sec:cortex_debug_def} Facilitates GDB-based debugging, enabling step-through execution and register-level inspection.
        \item \textbf{Raspberry Pi Pico:} Provides essential SDK management and project configuration utilities.
        \item \textbf{Pico-Debug-Dashboard:} A custom-made extension designed to unify how VS Code invokes the compiler toolchain and communicates with the hardware. A detailed breakdown of its features, design, and graphical user interface is provided in Section \ref{sec:pico_debug_dashboard} \cite{pico_debug_dashboard}.
    \end{itemize}
    
    \item[Connectivity and Drivers:] The software components that bridge the host machine with the microcontroller. This includes OpenOCD as the GDB-to-CMSIS-DAP server and USB/HID drivers (such as \texttt{libhidapi} and \texttt{libusb}) required to maintain communication with the debug probe.
\end{description}

All software components, including toolchain binaries, SDK releases, and debugging middleware, were configured according to their respective official documentation and established best practices and efforts were made to harmonise the configuration parameters across both the virtual machine and containerised environments. This is to ensure that any observed differences in performance or usability could be attributed strictly to the underlying virtualisation architecture.
\subsection{Virtual Machine Implementation}
Unlike containers, standard VM software supports USB passthrough natively, allowing the hypervisor to redirect the physical USB debug probe directly into the guest virtual environment. Additionally, because the VM runs a full desktop environment, students utilise the standard Visual Studio Code application, which enjoys full access to the official Microsoft Marketplace. This allows students to use the native Microsoft Serial Monitor extension as an integrated tool, whilst keeping Minicom available as a terminal-based serial console fallback.

\subsection{Containerisation Implementation}
As graphical interfaces are not natively supported in Docker containers, \texttt{code-server} is used in place of the original VS Code to provide a browser-based equivalent, ensuring a consistent user experience while maintaining the isolation and reproducibility of containerised environments \cite{code_server}. However, because \texttt{code-server} is limited to the Open VSX Registry, the Microsoft Serial Monitor extension is unavailable in this web-based environment. To maintain a consistent workflow, Minicom is used as a terminal-based alternative to provide consistent monitor support, with one-click status bar initialisation supported by the Pico-Debug-Dashboard extension.

Furthermore, container port forwarding, device mapping, and privilege controls are managed dynamically via a \texttt{docker-compose} configuration \cite{docker_compose_docs}. On native Linux host systems, Docker can access the physical USB bus directly by mounting device volumes (such as \texttt{/dev/bus/usb} and serial \texttt{/dev/ttyACM*} ports) under a privileged container state. However, on Windows host systems, physical USB devices are owned by the host OS and cannot be shared natively. To resolve this, the \texttt{usbipd} utility must be installed on the Windows host to redirect physical USB devices to the Windows Subsystem for Linux (WSL2) \cite{docker_wsl_docs, dorssel_usbipd_win}. \texttt{usbipd} functions by encapsulating USB traffic within TCP/IP packets, thereby overcoming the physical device isolation of the virtualised WSL2 kernel and allowing the container's volume mappings to mount the debug probe natively \cite{docker_wsl_docs}. To make \texttt{usbipd} easier to use, a specialised script is provided to help users automatically initialise and manage the tool.
\subsection{Pico-Debug-Dashboard Extension}
\label{sec:pico_debug_dashboard}
We developed a custom Visual Studio Code extension, the \texttt{Pico-Debug-Dashboard}, because the official CMake and Raspberry Pi extensions cluttered the IDE's status bar with too many inconvenient buttons, lacked support for the \texttt{pico-debug} second-core debugger, and frequently caused student configuration errors by requiring manual selection of compiler "kits" and build "variants" \cite{pico_getting_started, pico_debug, pico_debug_dashboard}. Written in TypeScript, our extension unifies these tools into a single, clean interface (illustrated in Figure \ref{fig:vscode_dashboard}). It automates project bootstrapping, build setups, and flashing via GDB or the official \texttt{picotool} utility, utilised to flash and query Raspberry Pi microcontrollers over USB \cite{picotool_github}, by executing raw, pre-configured terminal commands directly in the background. Crucially, the extension dynamically monitors the host's USB bus to detect whether a standard Pico, an official DebugProbe, or a virtual \texttt{pico-debug} probe is connected, displaying the real-time hardware status directly on the status bar.

\begin{figure}[ht]
    \centering
    \includegraphics[width=0.9\textwidth]{vscode_dashboard_ui} 
    \caption{The custom \texttt{Pico-Debug-Dashboard} interface within VS Code. The sidebar panel (top-left) manages project bootstrapping and configuration template syncing. The blue status bar (bottom-left) provides quick-action buttons for building, flashing, and serial console monitoring, whilst the telemetry status item (bottom-right) dynamically displays the connected hardware connection state.}
    \label{fig:vscode_dashboard}
\end{figure}

\section{Probe implementations}
\subsection{Second Pico Implementation}

The second board implementation utilises an additional RP-based microcontroller device as a dedicated hardware debugger. Now referred to as the probe, the device is flashed with the \texttt{debugprobe.uf2} firmware, turning it into a composite USB device that supports both the ARM CMSIS-DAP protocol for debugging and a Communication Device Class (CDC) interface for serial console monitoring. Physical connections between the probe and the target device are established via a 5-pin debug header, providing the necessary SWD signals and the bidirectional UART bridge needed for console output (see Section \ref{pico_wiring}). Once the connection between the probe and the host machine is established via USB, the debugging session can be initialised through VS Code or by manually initialising the GDB-to-OpenOCD-to-Probe pipeline through a terminal environment with commands\cite{pico_getting_started}:
\begin{lstlisting}[
    language=bash,
    caption={OpenOCD Initialisation Command},
    label={lst:openocd_init},
    frame=single,
    basicstyle=\ttfamily\small,
    breaklines=true,
    mathescape=false
]
$ sudo openocd -f interface/cmsis-dap.cfg -f target/rp2040.cfg -c "adapter speed 5000"
\end{lstlisting}

For an RP2350-based device, use \texttt{-f target/rp2350.cfg} instead.

Subsequently, the \texttt{GDB} client is invoked to link the debug symbols to the remote target:

\begin{lstlisting}[
    language=bash,
    caption={GDB Client Initialisation},
    label={lst:gdb_init},
    frame=single,
    basicstyle=\ttfamily\small,
    breaklines=true,
    mathescape=false
]
$ gdb-multiarch <project_name>.elf
(gdb) target remote localhost:3333
\end{lstlisting}
Once the debugging pipeline is properly initialised, the platform provides a robust environment for system diagnostics, supporting a comprehensive set of \texttt{GDB} and \texttt{OpenOCD} features for hardware-level control:
\begin{itemize}
\item \textbf{Instruction-Level Control:} Enables precise program flow management through \texttt{step-into}, \texttt{step-over}, and \texttt{continue} operations, allowing the developer to isolate logic errors at the assembly or source code level.
\item \textbf{Hardware and Software Breakpoints:} The platform automatically switches between hardware-assisted breakpoints for Flash memory and software breakpoints for RAM to facilitate non-intrusive execution halting.
\item \textbf{System State Inspection:} Grants deep visibility into the processor by exposing the Program Counter (\texttt{PC}), Stack Pointer (\texttt{SP}), and general-purpose registers.
\item \textbf{Real-Time Memory and Register Analysis:} Allows for the dynamic examination and modification of memory addresses and peripheral register maps (\texttt{SVD} integration).
\item \textbf{Stack Traceback and Backtracing:} Provides the ability to traverse the function call hierarchy, which is critical for identifying recursive errors, stack overflows, and improper task transitions within a multi-threaded \texttt{FreeRTOS} environment.
\end{itemize}

As a primary diagnostic interface, \texttt{Minicom} is integrated into the development platform to ensure reliable serial communication across all environments. The monitor session is initialised via a standardized task-based command, which automatically detects the active communication port (e.g., \texttt{/dev/ttyACM0} or \texttt{/dev/ttyACM1}) and defaults to a baud rate of 115200 bps. Within the \texttt{Minicom} terminal, standard keyboard shortcuts provide access to the control menu and session termination. This terminal-based approach provides a platform-agnostic method for capturing low-level system logs and \texttt{printf} output, functioning reliably even when graphical serial interfaces are inaccessible.

To start serial monitoring inside of VS Code when Minicom fails to start, the Microsoft's Serial Monitor extension is installed to provide intuitive UI/UX for serial communication, enabling developers to monitor \texttt{UART} or \texttt{USB} standard output directly within the IDE. The monitor session is initialised by navigating to the serial monitor tab resides on the bottom panel, then choosing the target communication port, which is usually \texttt{/dev/ttyACM0} or \texttt{/dev/ttyACM1}, and matching the baud rate to the microcontroller's configuration, typically 115200 bps. To send messages to the target, a message prompt is presented at the bottom for convenient input.

\subsection{Second core implementation \- pico-debug}

The approach where one core is used to debug another one was implemented and shared by majbthrd on Github. He developed two projects: "pico-debug" for RP2040 \cite{pico_debug}, and "pico2-debug" for RP2350 boards \cite{pico2_debug}. To setup the debugging session there are ready .uf2 files provided in the repositories. Pico-debug can be placed in RAM or cache memory of the board, and each option has its own .uf2 file. We are going with a more common and easier way to use it, placing it in the RAM, so the whole cache is available for the code of the program that is being debugged. To do so, we take the pico-debug-gimmecache.uf2 file from the releases page of the repository and flush it onto the board. When the application is designed to use the full RAM the pico-debug-maxram.uf2 is used instead.

Now the board shows as a CMSIS-DAP probe in the list of devices, and we can proceed to connect to it through the GDB and OpenOCD and flush just like we did with the second Pico board. The OpenOCD, however, has to be built by the user, since the newer versions of its Raspberry Pi's fork are not compatible with pico-debug. The detailed instructions on building OpenOCD are provided in the repository's "howto" section \cite{pico_debug} and involve cloning the repository and then compiling with make command. One big advantage is that this option does not need any additional hardware and cabels. 

\subsection{SEGGER J-link}

\begin{figure}[H]
    \begin{center}
      \includegraphics[width=0.48\textwidth]{jlink\_setup}
    \end{center}
    \caption{Physical setup of the J-Link implementation; the JTAG cable is connected to the J-Link EDU.}
    \label{fig:jlink_setup}
\end{figure}

The SEGGER J-Link, specifically the J-Link EDU V11 model, implementation requires a lot of background setup work before it becomes a viable debugging solution. Since the target device in this case the TKJHAT board that is already attached to a Single Board Computer (SBC) in the form of the Pico W, a direct connection via the debugging pins is possible. Wiring can be difficult without the necessary breakout boards, as loose cables can disrupt debugging operations. This specific implementation utilises male-to-male DuPont wire to bridge the capability gap from the lack of of proper proprietary connectors.

\begin{figure}[ht]
    \begin{center}
      \includegraphics[width=0.8\textwidth]{jlink\_connections}
    \end{center}
    \caption{List of connections from the J-Link to the Pico W. JTAG pinout diagram from OLIMEX ARM-JTAG dongle manual \cite{olimex}. Copyright~\textcopyright~2004, OLIMEX Ltd., All rights reserved.}
    \label{fig:jlink_connect}
\end{figure}


It is extremely important to correctly wire up this setup. Incorrect connections will most commonly result in a complete failure to establish communication between the SEGGER J-Link EDU V11 and the Pico W. If key signals such as SWDIO, SWCLK, or ground are miswired or not properly connected, the debugger will not be able to detect the target device at all. This typically appears as connection errors in the J-Link software, where the system repeatedly fails to identify or communicate with the microcontroller.

More severe consequences can occur if power-related pins are connected incorrectly. Miswiring voltage lines or shorting power to ground can damage the Pico W, the J-Link device, or both. Even less destructive errors, such as incorrect reference voltage connections, can lead to improper signal levels and unpredictable operation. For this reason, careful verification of all wiring is essential before attempting to use the debugger.

The Segger J-Link EDU probe requires initial firmware configuration and server setup before use. The firmware update is performed using the JLinkConfig utility, after which the JLinkGDBServer application provides the GDB server interface where it acts a bridge between GDB and the hardware. There exists a GUI for this tool, however, through testing, it has proven to be inflexible. In the instances where certain settings or physical connections are incorrect, the GUI will close automatically without giving the user an opportunity to try again. Therefore, for convenience sake, the command line is preferred for its supperior usability and reliability.

\begin{lstlisting}[
  language=bash,
  caption={J-Link EDU Command},
  label={lst:openocd_init},
  frame=single,
  basicstyle=\ttfamily\small,
  breaklines=true,
  mathescape=false
]
$ /opt/SEGGER/JLink/JLinkGDBServer select USB -device RP2040\_M0\_0 -endian little -if SWD -speed 4000 -noir -LocalhostOnly -nologtofile -port 2331 -SWOPort 2332 -TelnetPort 2333
\end{lstlisting}

The command will print some helpful messages about the process as it goes through the entire checklist of steps. The configuration requires selecting the target device as "RP2040\_M0\_0" and the target interface as "SWD". All remaining settings can be left at their default values. For security purposes, enabling the "Localhost only" option is recommended to restrict GDB server access to local connections only, preventing unauthorised remote access.

Unfortunately, the system demonstrated a high degree of unreliability during testing, primarily due to the quality of the physical connections between the target device and the debugging probe. The use of male-to-male DuPont wires on a breadboard introduced instability into the setup, making it difficult to maintain consistent electrical contact across the required signal lines. As a result, the SEGGER J-Link GDB Server repeatedly failed to detect the 3.3 volt reference signal from the Pico W, which is essential for establishing proper communication with the target. Without this reference voltage, the debugger is unable to determine the correct logic levels, effectively rendering the board undetectable from the software perspective.

This issue manifested in the debugging process as repeated connection failures and stalled initialisation attempts. Specifically, the GDB server would hang during the stage where it attempts to read and verify the target voltage, preventing any further progress in establishing a debugging session. In some instances, the process would terminate entirely after failing to complete this step, requiring manual intervention to restart the connection attempt. This behaviour significantly hindered the usability of the debugging setup and made consistent testing impractical.

A practical solution to this problem would involve improving the reliability and stability of the physical connections. One effective approach is the use of proper breakout boards or adapter modules that provide secure and well aligned pin interfaces for both the Pico W and the J-Link probe. These components reduce the likelihood of intermittent contact and ensure that critical signals, particularly ground and reference voltage, are consistently maintained. By replacing loose wiring with more robust connection methods, the overall stability of the debugging environment can be significantly improved, allowing the J-Link system to function as intended.

\section{3D-printed Casing}
Designing and producing a 3D-printed casing solution was paramount to ensuring that the SBC and HAT board hardware was integrated into a compact and durable debugging unit. The casing needed to be ergonomic, cheap and intuitive to use while providing the electronic components inside protection from damage that might occur during the usage of the  debugging unit. This requirement is reflected in the reinforced ridges along the chassis of the final casing that reduces the risk of cracking along the length of the case that might occur to accidental drops or abrasions.

Parametric design using computer-aided design (CAD) software was therefore an essential component of the development process, as it allowed the design to be systematically refined to meet the stringent functional and usability requirements associated with an effective educational debugging solution. By defining key dimensions and constraints as adjustable parameters, the design could be iteratively modified with precision, ensuring compatibility with the hardware while also accommodating practical considerations such as ease of assembly, durability, and user interaction. This approach enabled rapid evaluation of multiple design variations without requiring a complete redesign at each stage, which was particularly valuable given the need to balance technical performance with accessibility for student use.

The overarching design philosophy adopted throughout the project placed a strong emphasis on simplicity and user comfort, prioritising these factors above more complex or feature-heavy alternatives. This focus was driven by the intended educational context, where reliability, ease of use, and intuitive handling are more beneficial than intricate mechanical features that may introduce unnecessary points of failure. As a result, design decisions consistently favoured straightforward geometries, minimal part counts, and ergonomic considerations that would allow students to interact with the device confidently and without additional instruction.

\begin{figure}[H]
    \begin{center}
      \includegraphics[width=0.5\textwidth]{early_case}
    \end{center}
    \caption{Early version of the 3D-printed casing. Note the hinge-latch cover system.}
    \label{fig:debug_case_1}
\end{figure}

This philosophy is clearly reflected in the final iteration of the enclosure, particularly in the implementation of a simple inverted dovetail lid mechanism. Earlier versions of the design incorporated a more complex hinge and locking system, which, while functional, introduced additional manufacturing complexity and potential durability concerns. By replacing this mechanism with an inverted dovetail interface, the design achieves secure closure through a purely geometric solution that is both robust and easy to manufacture. This not only reduces the likelihood of mechanical failure but also simplifies assembly and everyday use, aligning with the broader objective of creating a practical and reliable debugging tool for students.

\begin{figure}[H]
  \begin{center}
    \includegraphics[width=0.5\textwidth]{case}
  \end{center}
  \caption{Version 6, with the inverted dovetail cover, of the 3D-printed casing. This version would be the final version of the case.}
  \label{fig:debug_case_2}
\end{figure}

The overall cumulative cost of raw materials required for manufacturing represented another significant and non-trivial consideration that was carefully evaluated throughout the iterative design process, given that there existed a clearly identified need to potentially scale the deployment of the debugging platform. After careful evaluation of available manufacturing materials, and due to the favorable combination of affordability characteristics and relative structural durability properties, Polylactic Acid (PLA) thermoplastic filament was selected and designated as the chosen material from which the protective enclosure would be fabricated and produced. The design was additionally and deliberately optimised specifically for ease of manufacturing using commercially available and widely accessible fused deposition modeling (FDM) three-dimensional printing equipment that is commonly present in many professional fabrication facilities, educational institutions, and even home-based maker and fabrication set-ups. 

Through careful iterative optimisation of the design geometry and thoughtful selection of print parameters, the final optimised design iteration achieved a remarkably efficient combined printing time of approximately two hours for the complete assembly and required only 26 grams of PLA filament material for the entire protective casing structure, metrics that clearly indicate and demonstrate a high level of successful optimization with regard to both print settings configuration and manufacturing best practices. The final version of the case also retained the modularity that was envisioned from the start, with it being able to accommodate for all the possible debugging options; from using the available second core present on the SBC to using another SBC to serve as the debugging probe.

\begin{figure}[H]
  \begin{center}
    \includegraphics[width=0.5\textwidth]{final_case}
  \end{center}
  \caption{The final version of the 3D-printed casing with the TKJHAT inside.}
  \label{fig:debug_case_3}
\end{figure}

\begin{figure}[H]
  \begin{center}
    \includegraphics[width=0.5\textwidth]{final_case_second_pico}
  \end{center}
  \caption{The final version of the 3D-printed casing with wiring set-up for a second pico debugger.}
  \label{fig:debug_case_4}
\end{figure}